\chapter{Complexity — P Subset P Poly Tableau}
%%% generated by the blueprint pipeline from TCSlib.Complexity.CircuitComplexity.PSubsetPPolyTableau
\section{Overview}

Declarations of the Lean module \texttt{TCSlib.Complexity.CircuitComplexity.PSubsetPPolyTableau}.

\section{Declarations}

\begin{definition}[Snapshot width]
\lean{Complexity.snapWidth}
\label{Complexity.snapWidth}
\leanok
\uses{Complexity.Snapshot}
For a finite Turing machine $M$ over the Boolean alphabet, its snapshot width is the
number of possible snapshots of $M$.
\end{definition}

\begin{definition}[Snapshot encoding]
\lean{Complexity.snapEncode}
\label{Complexity.snapEncode}
\leanok
\uses{Complexity.Snapshot, Complexity.snapWidth}
Let $M$ be a finite Turing machine over the Boolean alphabet and let $s$ be one of its
snapshots. Encode $s$ by a Boolean list of length equal to the number of possible
snapshots, with a $1$ at the position corresponding to $s$ and $0$ at every other
position.
\end{definition}

\begin{theorem}[Length of a snapshot encoding]
\lean{Complexity.length_snapEncode}
\label{Complexity.length_snapEncode}
\leanok
\uses{Complexity.Snapshot, Complexity.snapEncode, Complexity.snapWidth}
\difficulty{1}
Let $M$ be a finite-state Turing machine over the binary alphabet, and let $s$ be a
snapshot of $M$. The bit-string encoding of $s$ has length equal to the number of
possible snapshots of $M$.
\end{theorem}

\begin{theorem}[Injectivity of Snapshot Encoding]
\lean{Complexity.snapEncode_injective}
\label{Complexity.snapEncode_injective}
\leanok
\uses{Complexity.Snapshot, Complexity.snapEncode, Complexity.snapWidth}
\difficulty{4}
For every finite Turing machine $M$ over the Boolean alphabet, the encoding map from
snapshots of $M$ to finite Boolean strings is injective.
\end{theorem}

\begin{definition}[Snapshot decoder]
\lean{Complexity.snapDecode}
\label{Complexity.snapDecode}
\leanok
\uses{Complexity.Snapshot, Complexity.snapEncode}
Let $M$ be a finite Turing machine over the Boolean alphabet, and let $l$ be a finite
list of Boolean values. The snapshot decoded from $l$ is the value of the inverse of the
snapshot encoding of $M$ applied to $l$.
\end{definition}

\begin{theorem}[Decoding an Encoded Snapshot]
\lean{Complexity.snapDecode_snapEncode}
\label{Complexity.snapDecode_snapEncode}
\leanok
\uses{Complexity.Snapshot, Complexity.snapDecode, Complexity.snapEncode, Complexity.snapEncode_injective}
\difficulty{1}
Let $M$ be a finite Turing machine over the Boolean alphabet, and let $s$ be a snapshot
of $M$. Decoding the encoding of $s$ returns $s$.
\end{theorem}

\begin{definition}[Tableau arity]
\lean{Complexity.tableauArity}
\label{Complexity.tableauArity}
\leanok
\uses{Complexity.snapWidth}
For a finite Turing machine $M$ over the Boolean alphabet, define its tableau arity to
be $S+3+k(S+1)+1$, where $S$ is the number of possible snapshots of $M$ and $k$ is its
number of work tapes.
\end{definition}

\begin{definition}[Boolean tableau update]
\lean{Complexity.tableauStep}
\label{Complexity.tableauStep}
\leanok
\uses{Complexity.emitted, Complexity.snapDecode, Complexity.snapEncode, Complexity.snapWidth, Complexity.stepState, Complexity.writtenOrKept}
Let $M$ be a finite Turing machine over the Boolean alphabet, and let $l$ be any finite
Boolean list. Let $B$ be the width of a snapshot encoding; decode the first $B$ entries
of $l$ as a snapshot. Read the input symbol as the entry at position $B+2$ if the entry
at position $B+1$ is true, and otherwise take it to be absent. For each work tape, take
the chunk of $B+1$ entries beginning at position $B+3$ plus the tape index times $B+1$;
when its first entry is true, use the decoded tail as the scanned symbol after applying
the machine's write-or-keep operation, and otherwise take the scanned symbol to be
absent. If the entry at position $B$ is true, the output is the encoding of the initial
state, the read input symbol, and absent symbols on every work tape, followed by false.
Otherwise, the output is the encoding of the next state and the scanned symbols just
specified, together with the read input symbol, followed by the disjunction of the entry
at position $B+3+k(B+1)$ and whether the decoded snapshot emits true; entries beyond the
end of $l$ are read as false.
\end{definition}

\begin{definition}[Tableau update circuits]
\lean{Complexity.tableauGadgets}
\label{Complexity.tableauGadgets}
\leanok
\uses{BoolCircuit.DAGCircuit, BoolCircuit.gadget, Complexity.snapWidth, Complexity.tableauArity, Complexity.tableauStep}
Let $M$ be a finite Turing machine with $k$ work tapes. The tableau update circuits are
a list of $B+1$ Boolean circuits, where $B$ is the number of bits used to encode a
snapshot. Each circuit has input arity $B+3+k(B+1)+1$; the circuit at index $j$ outputs
bit $j$ of the one-step tableau update applied to the Boolean input vector.
\end{definition}

\begin{theorem}[Number of tableau gadgets]
\lean{Complexity.length_tableauGadgets}
\label{Complexity.length_tableauGadgets}
\leanok
\uses{Complexity.snapWidth, Complexity.tableauGadgets}
\difficulty{1}
For every finite-state Turing machine $M$ over the binary alphabet, the list of tableau
gadgets has length equal to the number of possible snapshots of $M$ plus one.
\end{theorem}

\begin{definition}[Tableau width]
\lean{Complexity.tableauWidth}
\label{Complexity.tableauWidth}
\leanok
\uses{BoolCircuit.embedWidth, Complexity.tableauGadgets}
For a Turing machine $M$ with Boolean-valued transition data, define its tableau width
to be the embedding width of the tableau gadgets associated with $M$.
\end{definition}

\begin{definition}[Tableau constant vertex]
\lean{Complexity.tableauConst}
\label{Complexity.tableauConst}
\leanok
Let $M$ be a finite Turing machine over the Boolean alphabet, let $n$ be a natural
number, and let $b$ be a Boolean value. The tableau constant vertex is $n+1$ if $b$ is
true, and $n$ if $b$ is false.
\end{definition}

\begin{definition}[Tableau block base]
\lean{Complexity.tableauBase}
\label{Complexity.tableauBase}
\leanok
\uses{Complexity.tableauWidth}
Let $M$ be a finite Turing machine over the Boolean alphabet, and let $n$ and $t$ be
natural numbers. The base index of tableau block $t$ is $n+2+tW$, where $W$ is the
tableau width of $M$.
\end{definition}

\begin{definition}[Tableau vertex index]
\lean{Complexity.tableauVertex}
\label{Complexity.tableauVertex}
\leanok
\uses{BoolCircuit.embedOffset, Complexity.tableauBase, Complexity.tableauGadgets}
Let $M$ be a finite Turing machine over the Boolean alphabet, and let $n$, $t$, and $j$
be natural numbers. The vertex index for bit $j$ of tableau block $t$ is the base index
$n+2+tW$ plus the embedding offset of output $j$ among $M$'s tableau update circuits,
where $W$ is the tableau width of $M$.
\end{definition}

\begin{definition}[Tableau block indices]
\lean{Complexity.tableauBlock}
\label{Complexity.tableauBlock}
\leanok
\uses{Complexity.snapWidth, Complexity.tableauVertex}
Let $M$ be a finite Turing machine over the Boolean alphabet, and let $n$ and $t$ be
natural numbers. The tableau block is the list of vertex indices, in order, for bits
$j=0,\ldots,B-1$, where $B$ is the number of bits used to encode a snapshot and each
index is the tableau vertex for bit $j$ in block $t$.
\end{definition}

\begin{definition}[Tableau tape sources]
\lean{Complexity.tableauTapeSources}
\label{Complexity.tableauTapeSources}
\leanok
\uses{Complexity.prevVisit, Complexity.snapWidth, Complexity.tableauBlock, Complexity.tableauConst}
Let $M$ be a finite Turing machine with Boolean alphabet, let $n,t\in\mathbb{N}$, and
let $\tau$ be a tape of $M$. If the previous visit time of tape $\tau$ before time $t$
is undefined, the source list has $\operatorname{snapWidth}(M)+1$ copies of the constant
vertex for $n$ with value false. If that previous visit time is $s$, the source list
consists of the constant vertex for $n$ with value true, followed by the tableau block
for $M,n,s$.
\end{definition}

\begin{definition}[Sources for a tableau step]
\lean{Complexity.tableauSources}
\label{Complexity.tableauSources}
\leanok
\uses{Complexity.inputPosAt, Complexity.snapWidth, Complexity.tableauBlock, Complexity.tableauConst, Complexity.tableauTapeSources, Complexity.tableauVertex}
Let $M$ be a finite Turing machine with Boolean input, and let $n,t\in\mathbb{N}$. The
source list is formed by concatenating: $w$ copies of the constant-false tableau vertex
if $t=0$, where $w$ is the snapshot width of $M$, and otherwise the tableau block for
$M$ on input length $n$ at time $t-1$; three entries recording whether $t=0$, whether
the input-head position at time $t$ lies between $1$ and $n$, and that position minus
one when it does (or the constant-false vertex otherwise); the tape-source lists for
each work tape; and, finally, the constant-false vertex if $t=0$, or the vertex at
position $w$ in the tableau at time $t-1$ otherwise.
\end{definition}

\begin{definition}[Tableau gate sequence]
\lean{Complexity.tableauGates}
\label{Complexity.tableauGates}
\leanok
\uses{BoolCircuit.DAGGate, BoolCircuit.constGate, BoolCircuit.embedAll, Complexity.tableauBase, Complexity.tableauGadgets, Complexity.tableauSources}
Let $M$ be a finite Turing machine with binary tape symbols, and let $n$ be a
nonnegative integer. Define a sequence of gate lists $(G_t)_{t\ge 0}$ by
$G_0=[\text{false constant gate},\text{true constant gate}]$ and, for each $t\ge 0$, by
appending to $G_t$ the gates obtained by embedding all tableau gadgets for $M$, with
source vertices determined by $n$ and $t$ and base index determined by $n$ and $t$, to
form $G_{t+1}$.
\end{definition}

\begin{theorem}[Length of the tableau gate list]
\lean{Complexity.length_tableauGates}
\label{Complexity.length_tableauGates}
\leanok
\uses{Complexity.tableauGates, Complexity.tableauWidth}
\difficulty{4}
For every finite-state Turing machine $M$ with Boolean input alphabet and every
$n,t\in\mathbb{N}$, the gate list at time $t$ in the tableau construction has length
$2+tW$, where $W$ is the total width of the tableau gadgets for $M$.
\end{theorem}

\begin{theorem}[Length of a Flat Map with Uniform Output Length]
\lean{Complexity.CfgTableau.length_flatMap_const}
\label{Complexity.CfgTableau.length_flatMap_const}
\leanok
\difficulty{4}
Let $l$ be a finite list and let $f$ map each element to a finite list. If every list
$f(a)$ has length $c$, then the concatenation of the lists $f(a)$ as $a$ ranges over $l$
has length $|l|c$.
\end{theorem}

\begin{theorem}[Length of a tableau block]
\lean{Complexity.length_tableauBlock}
\label{Complexity.length_tableauBlock}
\leanok
\uses{Complexity.snapWidth, Complexity.tableauBlock}
\difficulty{1}
Let $M$ be a finite-state Turing machine over the binary alphabet, and let $n,t$ be
natural numbers. The list of tableau vertex indices for the snapshot coordinates at time
$t$ has length equal to the number of possible snapshots of $M$.
\end{theorem}

\begin{theorem}[Length of the tape source list]
\lean{Complexity.length_tableauTapeSources}
\label{Complexity.length_tableauTapeSources}
\leanok
\uses{Complexity.snapWidth, Complexity.tableauTapeSources}
\difficulty{3}
For every finite-state Turing machine $M$ with Boolean input alphabet, all
$n,t\in\mathbb{N}$, and each work tape $\tau$, the source list for tape $\tau$ at time $t$
has length $B+1$, where $B$ is the number of bits in a snapshot.
\end{theorem}

\begin{theorem}[Length of the tableau source list]
\lean{Complexity.length_tableauSources}
\label{Complexity.length_tableauSources}
\leanok
\uses{Complexity.CfgTableau.length_flatMap_const, Complexity.length_tableauTapeSources, Complexity.snapWidth, Complexity.tableauArity, Complexity.tableauBlock, Complexity.tableauConst, Complexity.tableauSources}
\difficulty{4}
For every finite-state Turing machine $M$ with Boolean input alphabet and all
$n,t\in\mathbb{N}$, the source list for time $t$ has length equal to the tableau arity
$B+3+k(B+1)+1$, where $B$ is the number of bits in a snapshot and $k$ is the number of
work tapes.
\end{theorem}

\begin{theorem}[Tableau vertex lies within its time block]
\lean{Complexity.tableauVertex_lt}
\label{Complexity.tableauVertex_lt}
\leanok
\uses{BoolCircuit.embedOffset_lt, Complexity.snapWidth, Complexity.tableauBase, Complexity.tableauGadgets, Complexity.tableauVertex, Complexity.tableauWidth}
\difficulty{4}
Let $C$ be a directed acyclic circuit on $n$ inputs, let $t,n\in\mathbb{N}$, and let
$j\in\mathbb{N}$ satisfy $j\le B$, where $B$ is the number of bits in a machine
snapshot. The vertex assigned to index $j$ in time block $t$ is strictly less than the
base index of time block $t+1$.
\end{theorem}

\begin{theorem}[Monotonicity of the tableau base]
\lean{Complexity.tableauBase_mono}
\label{Complexity.tableauBase_mono}
\leanok
\uses{Complexity.tableauBase, Complexity.tableauWidth}
\difficulty{3}
Let $M$ be a finite-state Turing machine over the binary alphabet, let $n\in\mathbb{N}$,
and let $s,t\in\mathbb{N}$ satisfy $s\le t$. Then the tableau base at time $s$ is at
most the tableau base at time $t$.
\end{theorem}

\begin{theorem}[Tableau constants lie below the base]
\lean{Complexity.tableauConst_lt}
\label{Complexity.tableauConst_lt}
\leanok
\uses{Complexity.tableauBase, Complexity.tableauConst}
\difficulty{3}
For every finite-state Turing machine $M$ with Boolean input and output, all natural
numbers $n,t$, and every Boolean value $b$, the index assigned to the constant
$b$—namely $n$ for false and $n+1$ for true—is strictly less than the tableau base
$n+2+tW$, where $W$ is the tableau width determined by $M$.
\end{theorem}

\begin{theorem}[Tableau block lies below the later base]
\lean{Complexity.tableauBlock_lt}
\label{Complexity.tableauBlock_lt}
\leanok
\uses{Complexity.tableauBase, Complexity.tableauBase_mono, Complexity.tableauBlock, Complexity.tableauVertex_lt}
\difficulty{3}
Let $M$ be a finite-state Turing machine over the binary alphabet, let
$n,s,t\in\mathbb{N}$, and suppose $s<t$. Every vertex index in the tableau block at time
$s$ is strictly less than the tableau base at time $t$.
\end{theorem}

\begin{theorem}[A previous visit occurs earlier]
\lean{Complexity.prevVisit_lt}
\label{Complexity.prevVisit_lt}
\leanok
\uses{Complexity.prevVisit}
\difficulty{3}
Let $M$ be a finite binary Turing machine, and let $n,t,s$ be natural numbers and $\tau$
a work-tape index. If the last previous visit of tape $\tau$ at time $t$ on the all-zero
input of length $n$ is $s$, then $s<t$.
\end{theorem}

\begin{theorem}[Tableau sources precede the base]
\lean{Complexity.tableauSources_lt}
\label{Complexity.tableauSources_lt}
\leanok
\uses{Complexity.prevVisit_lt, Complexity.tableauBase, Complexity.tableauBase_mono, Complexity.tableauBlock_lt, Complexity.tableauConst, Complexity.tableauConst_lt, Complexity.tableauSources, Complexity.tableauTapeSources, Complexity.tableauVertex_lt}
\difficulty{5}
For every finite-state Turing machine $M$ with Boolean input and output and all natural
numbers $n,t$, every index in the source list used to construct the tableau at time $t$
is strictly less than the tableau base $n+2+tW$, where $W$ is the tableau width
determined by $M$.
\end{theorem}

\begin{theorem}[Acyclicity of tableau gates]
\lean{Complexity.gatesAcyclic_tableauGates}
\label{Complexity.gatesAcyclic_tableauGates}
\leanok
\uses{BoolCircuit.GatesAcyclic, BoolCircuit.GatesAcyclic.append, BoolCircuit.constGate, BoolCircuit.gatesAcyclic_embedAll, Complexity.length_tableauGates, Complexity.length_tableauSources, Complexity.tableauBase, Complexity.tableauGates, Complexity.tableauSources_lt}
\difficulty{4}
Let $M$ be a deterministic Turing machine with a finite number of work tapes and Boolean
input and output symbols, and let $n,t$ be natural numbers. In the gate list for the
tableau of $M$ on an input of length $n$ through time $t$, every argument of gate $i$
has index strictly less than $n+i$.
\end{theorem}

\begin{definition}[Tableau circuit]
\lean{Complexity.tableauCircuit}
\label{Complexity.tableauCircuit}
\leanok
\uses{BoolCircuit.DAGCircuit, Complexity.gatesAcyclic_tableauGates, Complexity.length_tableauGates, Complexity.snapWidth, Complexity.tableauBase, Complexity.tableauGates, Complexity.tableauVertex, Complexity.tableauVertex_lt}
Let $M$ be a finite Turing machine with a binary input alphabet, and let
$n,T\in\mathbb{N}$. The tableau circuit has $n$ input vertices, the gate list
constructed through time $T+1$, and output vertex equal to the tableau vertex for the
output bit at time $T$.
\end{definition}

\begin{theorem}[Size of the tableau circuit]
\lean{Complexity.tableauCircuit_size}
\label{Complexity.tableauCircuit_size}
\leanok
\uses{BoolCircuit.DAGCircuit.size, Complexity.length_tableauGates, Complexity.tableauCircuit, Complexity.tableauWidth}
\difficulty{3}
For every finite-state Turing machine $M$ with Boolean input and output, and all natural
numbers $n$ and $T$, the tableau circuit with input width $n$ and time parameter $T$ has
size $n+2+(T+1)W$, where $W$ is the tableau width determined by $M$.
\end{theorem}

\begin{theorem}[Two-input tableau circuit]
\lean{Complexity.tableauCircuit_isFaninTwo}
\label{Complexity.tableauCircuit_isFaninTwo}
\leanok
\uses{BoolCircuit.DAGCircuit.IsFaninTwo, BoolCircuit.constGate, BoolCircuit.faninTwo_embedAll, BoolCircuit.gadget_isFaninTwo, Complexity.tableauCircuit, Complexity.tableauGadgets, Complexity.tableauGates}
\difficulty{5}
Let $M$ be a finite-state Turing machine over the binary alphabet, and let
$n,T\in\mathbb{N}$. The tableau circuit for $M$ with parameters $n,T$ is well formed,
and each of its gates reads at most two input vertices.
\end{theorem}
